\scanpage{128}
de la conjecture analogue déjà faite pour la dim relative $1$.

Pour mieux appréhender cette deuxième partie \inserted{b)} de la conjecture, il faudrait examiner de plus près ce que deviennent, par composition avec $\mathcal E_L^{\bar L}\mathrel{\text{---}}\mathcal E_U^{\bar L}$, les sections de $\mathcal E_L^{\bar L}$ sur $\mathcal E_K$ associées à un $V$, et pour commencer, étudier le sous-groupe \inserted{$\Delta_V$} de $\pi_{U/K}^{\bar L}$ image de $I_V$ --- qui, on le devine, jouera un rôle analogue à un groupe-bout. Dans quelle mesure des $V$ distinctes définissent-elles des $\unsure{\Lambda_V|\Delta_V}$ distinctes ?

Pour étudier cette question, introduisons un modèle $X$ de $L/K$ contenant $U$, $X$ propre, les comp. irréd. de \inserted{$D=$} $X-U$ des diviseurs réguliers \inserted{$D_i$ $(i\in I)$} et $X-U$ à croisements normaux. Pour $x\in\operatorname{supp}D$, considérons le hensélisé \inserted{(non strict)} $\tilde O_x$ de $\underline{O}_{X,x}$, et l'image inverse $\unsure{\bar U_x|\tilde U_x}$ de $U$ dessus. On a
\margin{127}
